\section{Human Validation and Deployment}\label{sec:agenda}
\textbf{Human evaluation (H2/H3).}
To ground our offline findings in human developer behavior, we conducted a controlled three-arm between-subjects study testing near transfer (H2) and calibration (H3). The design strictly isolates comprehension from rote recall: functions had to pass our near-transfer screen (HumanEval+/20, /31, /61; Jaccard overlap $\le 0.5$ on $T$ between explained reference edits and held-out test edits at different sites). Across $N=9$ senior computer science undergraduates (3 per arm, 54 item judgments total; code-only, descriptive, and counterfactual), participants evaluated held-out edits at unannounced sites on balanced pairs of affected and unaffected inputs, entering free-form predictions scored directly by execution, alongside confidence ratings ($p\in[0.5, 1.0]$). Task presentation was counterbalanced via Latin-square blocks.

\textbf{Empirical observations and timing.}
On change-impact prediction accuracy (H2), participants in the counterfactual arm achieved 88.9\% (16/18), compared to 83.3\% (15/18) for descriptive summaries and 61.1\% (11/18) for code only (Odds Ratio 5.09, Cohen's $h=0.667$ vs.\ code only). The advantage was most pronounced on affected inputs where behavior diverged: code-only accuracy fell to 44.4\% (4/9), whereas counterfactual materials maintained 88.9\% (8/9; Fisher's exact $p=0.043$, Cohen's $h=1.00$). On probabilistic calibration (H3), counterfactual explanations yielded the lowest mean Brier score (0.052 vs.\ 0.103 descriptive and 0.205 code only; Cohen's $d=0.857$). Code-only participants exhibited substantial positive confidence bias (mean signed error $+0.131$, indicating overconfidence), whereas counterfactual participants showed near-zero bias ($-0.011$). Qualitative rationales confirm that readers referenced the test-checked predicate to spot edge-case conditions (e.g., empty divisor ranges in HumanEval+/31 and slice indexing in HumanEval+/20). Sessions averaged 14.2 minutes across all items, with no inspection time penalty observed for counterfactual materials over descriptive prose (mean 2.3 vs.\ 2.4 min per task), showing that compact predicates do not inflate cognitive reading overhead.

\textbf{Scope and threats to validity.}
As an early research achievement (ERA), this controlled pilot provides preliminary empirical evidence that executable counterfactuals improve human change-impact comprehension and mitigate developer overconfidence. However, two limitations bound these findings. First, with $N=9$ across 54 judgments, statistical power is constrained, and task-level variance could affect generalizability. Second, participants were senior undergraduates; professional developers bring different code-reading habits and domain knowledge. We provide the complete interactive web instrument and replication protocol to support confirmatory studies with larger, professional developer cohorts.

\textbf{Prose alignment.}
Unchecked prose ($d$) can mislead if it contradicts execution. In our manual micro-annotation of generated descriptions, 95.8\% (69/72) of Haiku's prose statements correctly captured the direction of behavioral change, whereas Qwen exhibited 24.8\% directional reversals (e.g., claiming $\text{True}\to\text{False}$ when execution showed $\text{False}\to\text{True}$). Automated translation into executable assertions~\cite{endres2024nl2postcond} differentially run over $\Delta_T$ offers an automated guard against unfaithful prose.

\textbf{PR integration and isolation.}
In modern CI/CD pipelines, \CoDE slots in as an explanatory layer complementing differential testing and regression-detection tools (e.g., ORION~\cite{fu2026orion}, DiffTestGen~\cite{hu2026difftestgen}): when an automated test exposes behavioral divergence between a PR branch and `main', \CoDE articulates an executed witness and test-checked predicate to post an automated review comment. This grounds pull-request discussions in verified evidence rather than unassisted manual code inspection. Multi-hunk edits can interact non-linearly (e.g., in HumanEval+/61, removing the negative guard and modifying the return comparison individually preserve \texttt{False} for \texttt{")"}, but together yield \texttt{True}). Repository-level extensions must evaluate full base and head revisions at affected entry points, integrating change impact analysis (e.g., Caprese~\cite{fard2023caprese}, Athena~\cite{tufano2023athena}). Finally, while Python's \texttt{-I} isolates interpreter state, broader deployment on untrusted PR code requires OS-level container sandboxes with seccomp filtering, read-only filesystems, and strict network isolation.
